\uuid{SUMp}
\titre{Fonction majorée/minorée}
\auteur{bodin ; hanani ; mzari}
\organisation{exo7}
\licence{CC BY-NC-SA 4.0}
\datecreate{2020-01-20}
\niveau{L1}
\module{Analyse}
\chapitre{Continuité, limite et étude de fonctions réelles}
\sousChapitre{Autre}
\difficulte{3}
\theme{Notion de fonctions, fonction majorée, minorée}
\exo7qcmid{lille-1.continuite.009}

\contenu{
\begin{question}
\qclassification{123.00}
\qtags{motcle=fonction majorée/minorée}
Soient $f,g : \Rr \to \Rr$. Quelles sont les assertions vraies ?
\begin{answers}
    \bad{Si $f$ est bornée et $g$ majorée alors $f-g$ est bornée.}

    \good{Si $f$ bornée et $g$ majorée alors $f-g$ est minorée.}

    \bad{Si $f$ et $g$ sont minorées, alors $f \times g$ est minorée.}

    \bad{Si $f$ et $g$ sont minorées, alors $|f \times g|$ est bornée.}
\end{answers}
\begin{explanations}
La somme de deux fonctions majorées (resp. minorées) est majorée (resp. minorée). Ce n'est pas le cas pour le produit : par exemple
$f(x) = -1$ est minorée, $g(x) = \exp(x)$ aussi, mais $f \times g (x) = -\exp(x)$ ne l'est pas.
\end{explanations}
\end{question}
}
